\documentclass[14pt,a4paper]{extarticle}

% --- XeLaTeX setup for Ukrainian ---
\usepackage{fontspec}
\usepackage{polyglossia}
\setmainlanguage{ukrainian}
\setotherlanguage{english}
\setmainfont{DejaVu Serif}
\setsansfont{DejaVu Sans}
\setmonofont{DejaVu Sans Mono}

% --- Math ---
\usepackage{amsmath, amssymb, amsthm}
\usepackage{mathtools}
\usepackage{unicode-math}

% --- Layout ---
\usepackage[left=2.5cm, right=1.5cm, top=2cm, bottom=2cm]{geometry}
\usepackage{setspace}
\onehalfspacing

\usepackage{enumitem}
\usepackage{booktabs}
\usepackage{array}
\usepackage{hyperref}
\usepackage{titlesec}
\usepackage{xcolor}

\hypersetup{
    colorlinks=true,
    linkcolor=blue!60!black,
    urlcolor=blue!60!black
}

\titleformat{\section}{\large\bfseries}{\thesection.}{0.5em}{}
\titleformat{\subsection}{\normalsize\bfseries}{\thesubsection.}{0.5em}{}

% --- Theorem environments ---
\theoremstyle{definition}
\newtheorem{definition}{Визначення}[section]
\newtheorem{remark}{Зауваження}[section]

\theoremstyle{plain}
\newtheorem{theorem}{Теорема}[section]

% --- Custom commands ---
\newcommand{\vset}[1]{\mathcal{#1}}
\newcommand{\Real}{\mathbb{R}}

% ================================================================
\begin{document}

\begin{center}
    {\large\bfseries ДИПЛОМНА РОБОТА}\\[4pt]
    {\normalsize на тему:}\\[6pt]
    {\large\bfseries «Розробка та програмна реалізація алгоритму оптимізації\\
    розподілу вантажів у транспортній мережі»}\\[12pt]
\end{center}

\vspace{1.5cm}

% ================================================================
\section{Математична постановка задачі}

\subsection{Формалізація транспортної мережі}

Транспортну мережу формалізуємо як \textbf{мультимодальний зважений орієнтований граф}:
\[
    G = (V,\; E,\; \Phi,\; \Psi),
\]
де:
\begin{itemize}[leftmargin=2em]
    \item $V = \{v_1, v_2, \ldots, v_n\}$ --- скінченна множина вузлів мережі;
    \item $E \subseteq V \times V$ --- множина орієнтованих ребер (доріг);
    \item $\Phi : V \to \vset{A}$ --- функція атрибутів вузлів;
    \item $\Psi : E \to \vset{B}$ --- функція атрибутів ребер.
\end{itemize}

\begin{definition}
\textbf{Вузол} $v_i \in V$ характеризується кортежем атрибутів:
\[
    \Phi(v_i) = \bigl(\mathrm{id}_i,\; \lambda_i,\; \mu_i,\; \delta_i^{\min},\; \delta_i^{\max},\; \kappa_i\bigr),
\]
де $\mathrm{id}_i$ --- унікальний ідентифікатор; $(\lambda_i, \mu_i) \in \Real^2$ --- географічні координати (широта, довгота); $[\delta_i^{\min}, \delta_i^{\max}]$ --- часове вікно обслуговування; $\kappa_i \in \Real_{\geq 0}$ --- вантажопотужність терміналу.
\end{definition}

\begin{definition}
\textbf{Ребро} $e_{ij} = (v_i, v_j) \in E$ характеризується кортежем:
\[
    \Psi(e_{ij}) = \bigl(\ell_{ij},\; M_{ij}\bigr),
\]
де $\ell_{ij} \in \Real_{>0}$ --- довжина сегменту (метри); $M_{ij} \subseteq \mathcal{T}$ --- множина допустимих для цього ребра типів транспортних засобів.
\end{definition}

\subsection{Типи транспортних засобів та мультимодальність}

Визначимо множину типів транспортних засобів:
\[
    \mathcal{T} = \{\mathrm{Car},\; \mathrm{Bike},\; \mathrm{Walk}\},
\]
та визначимо підмножину ребер, доступних для кожного типу транспорту:
\[
    E_\theta = \bigl\{\, e \in E \mid \theta \in M_e \,\bigr\}, \qquad \theta \in \mathcal{T}.
\]
Підмножини $E_\theta$ можуть перетинатися: одне ребро може бути доступним для кількох типів транспорту.

Для кожного типу транспорту $\theta \in \mathcal{T}$ введемо характеристичну швидкість $\bar{u}_\theta$ (м/с) та вантажну місткість одиниці $Q_\theta$ (кг):
\[
    \bar{u}_{\mathrm{Car}} > \bar{u}_{\mathrm{Bike}} > \bar{u}_{\mathrm{Walk}}, \qquad
    Q_{\mathrm{Car}} \geq Q_{\mathrm{Bike}} \geq Q_{\mathrm{Walk}}.
\]

\textbf{Парк транспортних засобів.} Визначимо скінченну множину наявних транспортних засобів:
\[
    \mathcal{V} = \{v_1^{\mathcal{V}}, v_2^{\mathcal{V}}, \ldots, v_F^{\mathcal{V}}\}, \qquad F = |\mathcal{V}|,
\]
де кожному транспортному засобу $w \in \mathcal{V}$ відповідає тип $\theta_w \in \mathcal{T}$. Кількість одиниць кожного типу:
\[
    F_\theta = \bigl|\{\, w \in \mathcal{V} \mid \theta_w = \theta \,\}\bigr|, \qquad \theta \in \mathcal{T}.
\]
Вантажна місткість транспортного засобу $w$ дорівнює $Q_{\theta_w}$.

\subsection{Перевірка доступності транспорту для проходження ребра}

Початок маршруту та кінцева точка доставки можуть бути поза межами графу $G$, тому
введемо метрику доступності типу транспорту $\theta$ від вузла $v \in V$ як відстань до найближчого вузла, з якого виходить хоча б одне ребро типу $\theta$:
\[
    d\bigl(v,\, E_\theta\bigr) = \min_{\substack{u \in V \\ \exists\, e \in E_\theta:\; v_i(e) = u}} \mathrm{dist}_G(v,\, u) \;\in\; \Real_{\geq 0},
\]
де $\mathrm{dist}_G(v, u)$ --- довжина найкоротшого шляху між $v$ та $u$ у графі $G$; $v_i(e)$ --- початковий вузол ребра $e$.

Тоді \textbf{логічна функція доступності} транспорту типу $\theta$ у вузлі $v$ визначається як:
\[
    \mathcal{F}_\theta(v) =
    \begin{cases}
        1, & \text{якщо } d\bigl(v,\, E_\theta\bigr) < 100 \text{ м},\\
        0, & \text{інакше}.
    \end{cases}
\]

\textbf{Ієрархічна стратегія вибору транспорту} (пріоритет Car $\succ$ Bike $\succ$ Walk) формалізується через систему вкладених умов:

\[
    \theta^*(v) =
    \begin{cases}
        \mathrm{Car},  & \text{якщо } \mathcal{F}_{\mathrm{Car}}(v) = 1, \\[4pt]
        \mathrm{Bike}, & \text{якщо } \mathcal{F}_{\mathrm{Car}}(v) = 0 \;\wedge\;
                                      \mathcal{F}_{\mathrm{Bike}}(v) = 1, \\[4pt]
        \mathrm{Walk}, & \text{якщо } \mathcal{F}_{\mathrm{Car}}(v) = 0 \;\wedge\;
                                      \mathcal{F}_{\mathrm{Bike}}(v) = 0 \;\wedge\;
                                      \mathcal{F}_{\mathrm{Walk}}(v) = 1, \\[4pt]
        \varnothing,   & \text{інакше (вузол недосяжний).}
    \end{cases}
\]

\begin{remark}
Умова $d(v, E_\theta) < 100$ м було обрано як відстань, яку кур'єр може пронести вантаж пішки і може бути параметризована відповідно до конкретної мережі.
\end{remark}

\subsection{Формалізація задачі маршрутизації}

Нехай задано множину вантажів $\mathcal{K} = \{k_1, k_2, \ldots, k_m\}$. Кожен вантаж $k \in \mathcal{K}$ характеризується кортежем:
\[
    k = \bigl(o_k,\; d_k,\; q_k\bigr),
\]
де $o_k, d_k \in V$ --- вузли відправлення та призначення; $q_k \in \Real_{>0}$ --- маса вантажу.

Введемо \textbf{функцію призначення} $\xi : \mathcal{K} \to \mathcal{V}$, яка кожному вантажу ставить у відповідність транспортний засіб, що його доставляє. Множина вантажів, призначених транспортному засобу $w \in \mathcal{V}$:
\[
    \mathcal{K}_w = \bigl\{\, k \in \mathcal{K} \mid \xi(k) = w \,\bigr\}.
\]

\textbf{Маршрутом} для вантажу $k$ назвемо послідовність:
\[
    P_k = \bigl(v_0^{(k)}, e_1^{(k)}, v_1^{(k)}, e_2^{(k)}, \ldots, e_{r_k}^{(k)}, v_{r_k}^{(k)}\bigr),
\]
де $v_0^{(k)} = o_k$, $\; v_{r_k}^{(k)} = d_k$, \; $e_j^{(k)} \in E$; $r_k$ --- кількість ребер маршруту $P_k$.

\textbf{Час проходження} сегменту $e_{ij}$ транспортним засобом $\theta^*(v_i)$:
\[
    t_{ij}^{\bigl(\theta^*(v_i)\bigr)} = \frac{\ell_{ij}}{\bar{u}_{\theta^*(v_i)}} + \tau_{ij}^{\text{перес}},
\]
де $\tau_{ij}^{\text{перес}} \geq 0$ --- час пересадки при зміні виду транспорту на ребрі $e_{ij}$ (дорівнює нулю, якщо вид транспорту не змінюється).

\subsection{Цільова функція}

\textbf{Сумарний час доставки} по всій множині вантажів визначається як:
\[
    T\bigl(\{P_k\}_{k \in \mathcal{K}}\bigr)
    = \sum_{k \in \mathcal{K}} T_k(P_k)
    = \sum_{k \in \mathcal{K}} \sum_{j=1}^{r_k} t_{e_j^{(k)}}^{\bigl(\theta^*(v_{j-1}^{(k)})\bigr)},
\]
де $T_k(P_k)$ --- сумарний час доставки вантажу $k$ маршрутом $P_k$.

\textbf{Задача оптимізації} полягає у знаходженні набору маршрутів $\{P_k^*\}_{k \in \mathcal{K}}$ та функції призначення $\xi^*$, які мінімізують цільову функцію:

\[
    \boxed{
    \min_{\{P_k\},\, \xi} T\bigl(\{P_k\}\bigr)
    = \min_{\{P_k\},\, \xi} \sum_{k \in \mathcal{K}} \sum_{j=1}^{r_k} t_{e_j^{(k)}}^{\bigl(\theta^*(v_{j-1}^{(k)})\bigr)}
    }
\]

за умови виконання таких обмежень:

\begin{enumerate}[label=\textbf{О\arabic*.}, leftmargin=3em]
    \item \textbf{(Зв'язність маршруту):} $\forall k \in \mathcal{K}$, $\forall j \in \{1,\ldots,r_k\!-\!1\}$: кінцевий вузол $e_j^{(k)}$ є початковим вузлом $e_{j+1}^{(k)}$.
    \item \textbf{(Вантажна місткість):} $\forall w \in \mathcal{V}$: $\displaystyle\sum_{k \in \mathcal{K}_w} q_k \leq Q_{\theta_w}$.
    \item \textbf{(Часові вікна):} $\forall k \in \mathcal{K}$: $a_k \leq T_k(P_k) \leq b_k$.
    \item \textbf{(Доступність транспорту):} $\forall k \in \mathcal{K}$, $\forall v \in P_k$: $\theta^*(v) \neq \varnothing$.
    \item \textbf{(Допустимість типу транспорту):} $\forall e_{ij} \in P_k$: $\theta^*(v_i) \in M_{ij}$.
    \item \textbf{(Узгодженість засобу та маршруту):} $\forall k \in \mathcal{K}$, $\forall e_{ij} \in P_k$: $\theta_{\xi(k)} \in M_{ij}$.
    \item \textbf{(Цілісність маршруту):} $v_0^{(k)} = o_k$ та $v_{r_k}^{(k)} = d_k$, $\; \forall k \in \mathcal{K}$.
\end{enumerate}

\end{document}
